\chapter{Postgres desde o inicio}

\section{Por que Postgres em vez de CSV?}

CSV e simples para debugar, mas o objetivo deste projeto tambem e treinar banco, indices, queries e performance. Entao faz sentido usar Postgres desde o inicio.

Com Postgres voce treina:

\begin{itemize}
  \item schema design;
  \item tipos de dados;
  \item \code{JSONB};
  \item batch insert;
  \item transacoes;
  \item indices;
  \item \code{EXPLAIN ANALYZE};
  \item queries de insight.
\end{itemize}

\section{Docker Compose para Postgres}

Arquivo: \code{docker-compose.yml}

\begin{lstlisting}[style=devbutterbash]
services:
  postgres:
    image: postgres:16
    container_name: apple_health_postgres
    environment:
      POSTGRES_USER: apple
      POSTGRES_PASSWORD: apple
      POSTGRES_DB: apple_health
    ports:
      - "5433:5432"
    volumes:
      - postgres_data:/var/lib/postgresql/data

volumes:
  postgres_data:
\end{lstlisting}

Subir banco:

\begin{lstlisting}[style=devbutterbash]
docker compose up -d
\end{lstlisting}

\section{Settings}

Arquivo: \code{src/config/settings.py}

\begin{lstlisting}[style=devbutterpython]
from dataclasses import dataclass
import os


@dataclass(frozen=True, slots=True)
class Settings:
    postgres_host: str = os.getenv("POSTGRES_HOST", "localhost")
    postgres_port: int = int(os.getenv("POSTGRES_PORT", "5433"))
    postgres_user: str = os.getenv("POSTGRES_USER", "apple")
    postgres_password: str = os.getenv("POSTGRES_PASSWORD", "apple")
    postgres_db: str = os.getenv("POSTGRES_DB", "apple_health")
    batch_size: int = int(os.getenv("BATCH_SIZE", "5000"))

    @property
    def database_url(self) -> str:
        return (
            f"postgresql://{self.postgres_user}:"
            f"{self.postgres_password}@"
            f"{self.postgres_host}:{self.postgres_port}/"
            f"{self.postgres_db}"
        )
\end{lstlisting}

\section{Schema inicial}

Tabela principal: \code{health_records}.

\begin{lstlisting}[style=devbuttersql]
CREATE TABLE IF NOT EXISTS health_records (
    id BIGSERIAL PRIMARY KEY,
    apple_record_id BIGINT NOT NULL,
    type TEXT,
    source_name TEXT,
    source_version TEXT,
    device TEXT,
    unit TEXT,
    creation_date TIMESTAMPTZ,
    start_date TIMESTAMPTZ,
    end_date TIMESTAMPTZ,
    value TEXT,
    value_numeric DOUBLE PRECISION,
    metadata JSONB
);
\end{lstlisting}

\section{Tabela parse\_runs}

Essa tabela guarda historico de execucoes do parser.

\begin{lstlisting}[style=devbuttersql]
CREATE TABLE IF NOT EXISTS parse_runs (
    id BIGSERIAL PRIMARY KEY,
    started_at TIMESTAMPTZ NOT NULL DEFAULT NOW(),
    finished_at TIMESTAMPTZ,
    file_path TEXT NOT NULL,
    total_records BIGINT DEFAULT 0,
    status TEXT NOT NULL DEFAULT 'running',
    error TEXT
);
\end{lstlisting}

\section{Indices}

Consultas de insights vao filtrar muito por tipo e data. Entao os primeiros indices sao:

\begin{lstlisting}[style=devbuttersql]
CREATE INDEX IF NOT EXISTS idx_health_records_type
ON health_records (type);

CREATE INDEX IF NOT EXISTS idx_health_records_start_date
ON health_records (start_date);

CREATE INDEX IF NOT EXISTS idx_health_records_type_start_date
ON health_records (type, start_date);
\end{lstlisting}

\section{Por que indice importa?}

Sem indice, o banco pode precisar varrer muitas linhas:

\[
O(n)
\]

Com indice B-tree, buscas por tipo/data podem se aproximar de:

\[
O(log n)
\]

Na pratica, confirme com:

\begin{lstlisting}[style=devbuttersql]
EXPLAIN ANALYZE
SELECT
    DATE(start_date) AS day,
    AVG(value_numeric) AS avg_heart_rate
FROM health_records
WHERE type = 'HKQuantityTypeIdentifierHeartRate'
GROUP BY DATE(start_date)
ORDER BY day;
\end{lstlisting}
